\begin{abstract}
\noindent Una colocación de reinas tiene a lo sumo una reina en cada fila, columna y diagonal de ambas direcciones.
Estudiamos el problema continuo de maximizar la entropía de una densidad de masa
$\rho$ en el cuadrado unidad bajo estas cuatro restricciones. Para $\rho=1$ el funcional es el del problema variacional
de Simkin para las $n$ reinas. Para $0<\rho<1$, el maximizador es único
y su densidad en cada punto, salvo en un conjunto de medida cero, es proporcional al producto de
las capacidades restantes de las cuatro líneas que pasan por él; la constante de proporcionalidad
$\lambda(\rho)$ es la actividad.
Al acercarse la masa a uno, la capacidad restante de filas y columnas tiende a cero.
Al dividirla por $x=1-\rho$ obtenemos perfiles positivos que convergen uniformemente
y, junto con las capacidades restantes de las diagonales y el límite de la actividad reescalada $\lambda x^2$, determinan el maximizador de masa uno. Estos perfiles dependen analíticamente de $x$ cerca de cero
y satisfacen una cota Lipschitz espacial uniforme hasta el tablero lleno.
La actividad $\lambda(\rho)$ del maximizador determina la función
$F(x)=(1-x)/(\lambda(1-x)x^2)$, que se extiende analíticamente a ambos extremos y cumple
$\int_0^1\log F(x)\,dx=1+s^*(1)$, donde $s^*(1)$ es la entropía óptima en masa uno; además,
$s^*(1-x)-s^*(1)=2x\log(1/x)+O(x)$. Una extrapolación numérica, sin certificación de error,
da $F(0)\approx0{,}0809$ y $s^*(1)\approx-1{,}9440$, compatible con los valores conocidos de la
constante de Simkin. Todos los resultados se refieren al problema continuo.
\end{abstract}
\noindent\small\textbf{Palabras clave.} Problema de las $n$ reinas, entropía variacional,
ecuaciones de Euler, regularidad espacial. \quad\textbf{MSC 2020.} 49K27, 05A16, 46N10, 94A17.\par\normalsize